\documentclass[11pt]{article}
\usepackage[T1]{fontenc}
\usepackage{lmodern}
\usepackage[margin=1in]{geometry}
\usepackage{amsmath,amssymb,amsthm}
\usepackage{microtype}
\usepackage[dvipsnames]{xcolor}
\usepackage[colorlinks=true,linkcolor=MidnightBlue,citecolor=MidnightBlue,urlcolor=MidnightBlue]{hyperref}
\usepackage{fancyhdr}
\pagestyle{fancy}
\fancyhf{}
\fancyhead[L]{\small SUSY positivity research program}
\fancyhead[R]{\small Mathematical background}
\fancyfoot[C]{\thepage}
\setlength{\headheight}{14pt}
\setlength{\emergencystretch}{2em}
\numberwithin{equation}{section}
\newtheorem{proposition}{Proposition}[section]
\hypersetup{pdftitle={SUSY positivity research program: mathematical background},pdfauthor={Working research document prepared for Edward Baker},pdfsubject={Shared Weil positivity framework, shifted transfers, and proposed bulk-boundary realizations},pdfkeywords={Weil positivity, Riemann hypothesis, shifted zeta, positive factorizations, research program}}
\begin{document}
\begin{center}
{\Large\bfseries SUSY positivity research program}\\[4pt]
{\large Mathematical background}\\[6pt]
{\small Prepared 11 September 2026; expanded 13 September 2026.}
\end{center}
\noindent This standalone background includes the definitions and
derivations needed for its positive kinetic construction and its
conditional routes to full Weil positivity. The full arithmetic matching remains an open construction problem.
\medskip
% Shared mathematical background for the SUSY positivity research program.
% Authoritative substantive section; background.tex supplies the standalone wrapper.
% Uses amsmath, amssymb, amsthm; the proposition environment; and bibliography
% keys SuzukiScrewBackground, DLMF and KwasnickiMuchaBackground,
% supplied in background.tex.

\section{Weil positivity, support length, and the shifted family}
\label{sec:rh-background}

This document develops the common mathematical framework of the
SUSY positivity research program. It gives the arithmetic target, the
shift-energy identities and their domains, and an explicit positive
auxiliary-field realization of the gamma kinetic term. Definitions and
derivations used within this framework are included here. External
references supply the classical Weil criterion, standard special-function
identities, and the general theory of positive extensions. A positive
realization of the full arithmetic form at arbitrary support length
remains the objective, not a conclusion of the constructions below.

The positivity problem has two parameters with different
roles. The length $L$ specifies the support available to a test function.
The shift $\omega$ belongs to an auxiliary family of zeta transfers and
their generators. The classical equivalence with the Riemann hypothesis
(RH) concerns the central form, at $\omega=0$, on inputs of arbitrary
compact support. We first state that target, then explain why shifted
forms and contraction identities are useful in seeking a positive
realization of it. We also distinguish a direct amplitude factor, a
bulk energy reduced to boundary data, and an energy accumulated along
the shift family.

\subsection{The arithmetic form and its input space}
\label{subsec:rh-form}

Let $I_L=(-L/2,L/2)$, so that $L>0$ is the total interval length.
All norms and inner products below use unweighted Lebesgue measure, with
inner products conjugate-linear in the first variable. Extend
$f\in L^2(I_L)$ by zero to $\widetilde f$ on $\mathbb R$, and set
\begin{equation}
\widehat f(\tau)=\int_{I_L}f(x)e^{-i\tau x}\,dx,\qquad
\mathcal D_{\log,L}=
\left\{f\in L^2(I_L):
\int_{\mathbb R}\log(2+|\tau|)|\widehat f(\tau)|^2\,d\tau<\infty\right\}.
\label{eq:rh-domain}
\end{equation}
For $a>0$ write $(T_af)(x)=\widetilde f(x-a)$ on $I_L$, and define
\[
C(f)=\int_{I_L}f(x)\cosh(x/2)\,dx,\qquad
S(f)=\int_{I_L}f(x)\sinh(x/2)\,dx.
\]
The central finite-interval Weil form in our normalization is
\begin{align}
Q_{0,L}[f]={}&
\frac1{2\pi}\int_{\mathbb R}
\left(\operatorname{Re}\psi\left(\tfrac14+\tfrac{i\tau}{2}\right)
-\log\pi\right)|\widehat f(\tau)|^2\,d\tau \notag\\
&+2|C(f)|^2-2|S(f)|^2
-\sum_{\substack{n\ge2\\\log n<L}}\frac{\Lambda(n)}{\sqrt n}
\langle f,(T_{\log n}+T_{\log n}^*)f\rangle .
\label{eq:rh-central}
\end{align}
Here $\psi=\Gamma'/\Gamma$ and $\Lambda$ is the von Mangoldt function.
The sum is finite and only prime powers contribute. A delay of length
$L$ has zero action, which explains the strict cutoff.

The multiplier in \eqref{eq:rh-central} differs from
$\log(2+|\tau|)$ by a bounded function, by the digamma asymptotics
\cite[\S5.11]{DLMF}. The remaining terms are bounded forms at each
fixed $L$. Consequently $Q_{0,L}$ is closed and semibounded on
$\mathcal D_{\log,L}$ without any positivity assumption.
Smooth compactly supported inputs form a core: dilation can first move
the support into the interval interior, and mollification then approximates
in the weighted Fourier norm in \eqref{eq:rh-domain}.

We write $Q^\gamma_{0,L}$ for the form obtained by omitting the explicit
prime sum. This convention retains the normalization and pole terms
$2|C|^2-2|S|^2$. It gives $Q^\gamma_{0,L}=Q_{0,L}$ for
$L\leq\log2$; beyond that range gamma positivity alone does not establish
the required arithmetic positivity.

\paragraph{Separating the positive kinetic term from the full form.}
The gamma multiplier has a frequency-dependent part and a constant part.
Define, for $s\geq0$,
\[
B(s)=\operatorname{Re}\psi\!\left(\tfrac14+\tfrac{i\sqrt{s}}2\right)
-\psi(\tfrac14),\qquad w_0=\psi(\tfrac14)-\log\pi.
\]
The constant is negative:
$w_0=-\gamma_E-\pi/2-3\log2-\log\pi<0$, where $\gamma_E$ is Euler's
constant. For the zero-extended input, set
\[
K[\widetilde f]
=\frac1{2\pi}\int_{\mathbb R}B(\tau^2)|\widehat f(\tau)|^2\,d\tau,
\qquad P[f]=2|C(f)|^2-2|S(f)|^2.
\]
The letter $B$ denotes a scalar frequency response; $K$ denotes the
quadratic functional obtained by weighting all Fourier components with
that response. The argument $s=\tau^2$ is squared spatial frequency,
not a time or a zeta shift.

Write $\ell_p=\log p$ and $r_p=p^{-1/2}$ for a prime $p$.
Since $\Lambda(p^m)=\log p$ and $\Lambda(n)=0$ off prime powers,
Plancherel's identity rewrites the full central form exactly as
\[
\boxed{\begin{aligned}
Q_{0,L}[f]={}&K[\widetilde f]+w_0\lVert f\rVert^2+P[f]\\
&-\sum_{\substack{p\ {\rm prime},\ m\geq1\\m\ell_p<L}}
\ell_pr_p^m\,
\langle f,(T_{m\ell_p}+T_{m\ell_p}^*)f\rangle .
\end{aligned}}
\]
In particular,
\[
Q^\gamma_{0,L}[f]=K[\widetilde f]+w_0\lVert f\rVert^2+P[f].
\]
Thus the gamma kinetic energy is $K$, while the gamma form in this
document includes the contact and pole terms as well. Even on a
prime-free interval, the full form is generally not equal to $K$.
This decomposition is an arithmetic identity; it does not assert
positivity of its signed terms or of their sum.

The positive structure of $B$ follows from the digamma partial-fraction
identity \cite[\S5.7]{DLMF}. With $a_k=2k+1/2$,
\[
\begin{aligned}
B(\tau^2)
&=\sum_{k=0}^\infty
\left(\frac{1}{k+1/4}
-\frac{k+1/4}{(k+1/4)^2+\tau^2/4}\right)\\
&=\sum_{k=0}^\infty\frac{2}{a_k}
\frac{\tau^2}{a_k^2+\tau^2}\geq0.
\end{aligned}
\]
Each difference is nonnegative, and the series converges with terms
$O_\tau(k^{-3})$. Therefore $K\geq0$ on its form domain. The later
bulk--boundary construction derives each summand as the minimum of a
positive field energy. The other terms in the boxed decomposition
specify what a full positive realization must still reproduce.

\subsection{The quantifiers in the RH criterion}
\label{subsec:rh-criterion}

Weil's positivity criterion, in its smooth compact-support formulation
\cite[\S1.1]{SuzukiScrewBackground}, states
\begin{equation}
\boxed{\mathrm{RH}\quad\Longleftrightarrow\quad
Q_{0,L}[f]\geq0
\quad\text{for every }L>0
\text{ and every }f\in C_c^\infty(I_L).}
\label{eq:rh-criterion}
\end{equation}
Complex-valued test functions are allowed. Every compactly supported
smooth function on $\mathbb R$ belongs to $C_c^\infty(I_L)$ for some
$L$, and \eqref{eq:rh-central} agrees on that input at every larger
length. Conversely, every such finite-interval input is a whole-line
test function. This explains the equivalence between the usual global
criterion and \eqref{eq:rh-criterion}. By the core property, one may
equivalently require nonnegativity on every $\mathcal D_{\log,L}$.

In particular, the first subscript in $Q_{0,L}$ has already fixed
$\omega=0$. The remaining quantifiers range over the length and
\emph{all} inputs, including inputs outside any chosen finite-dimensional
approximation. Positivity here means nonnegativity. A lower bound
$Q_{0,L}[f]\geq m_L\lVert f\rVert^2$ with $m_L>0$ is useful for
quantitative arguments, but the criterion requires no positive constant
independent of $L$.

An unbounded sequence of covered lengths suffices: any fixed input
eventually fits in one of them. A finite collection of lengths does not
suffice, and decomposing a long input into short pieces introduces cross
terms that must also be controlled.

\subsection{The shifted transfer and its generator}
\label{subsec:rh-shift}

Define the completed zeta function and its central translate by
\[
\xi(s)=\tfrac12s(s-1)\pi^{-s/2}\Gamma(s/2)\zeta(s),
\qquad F(p)=\xi(\tfrac12+p).
\]
For $0<\omega\leq1/2$ the causal transfer has Laplace symbol
\begin{equation}
K_\omega(p)=\frac{F(p-\omega)}{F(p+\omega)},\qquad
\Theta_\omega(z)=K_\omega(-iz)
=\frac{\xi(1/2-\omega-iz)}{\xi(1/2+\omega-iz)}.
\label{eq:rh-transfer}
\end{equation}
The causal realization is defined on a line $\operatorname{Re}p=\eta>1$,
where the relevant zeta Dirichlet series converge absolutely and the
denominator is nonzero. Let $k_\omega$ be its inverse Laplace kernel,
supported in $[0,\infty)$. Then
$V_{\omega,L}f=(k_\omega*\widetilde f)|_{I_L}$.
Equivalently, identify interval inputs with their zero extensions,
let $P_L$ multiply by $\mathbf1_{I_L}$ on the whole line, and let
$M_\eta$ multiply by $e^{-\eta x}$. Then
\[
V_{\omega,L}
=M_\eta^{-1}P_L\,\mathcal M[K_\omega(\eta+i\,\cdot)]\,P_LM_\eta ,
\]
where $\mathcal M[b]$ denotes the whole-line Fourier multiplier by $b$.
This defines a bounded finite-interval operator; it does not assume a
bounded whole-line transfer on the unweighted space. At zero shift
$K_0=1$ and $V_{0,L}=I$ as a strong limit.

The symmetric generator form can be specified directly on
$\mathcal D_{\log,L}$. Put
$R(t)=2\cosh(t/2)-e^{-t/2}/(1-e^{-2t})$ for $t>0$. Then
\begin{align}
Q_{\omega,L}[f]={}&Q_{0,L}[f]
+\int_{I_L^2}\bigl(\cosh(\omega|x-y|)-1\bigr)R(|x-y|)
\overline{f(x)}f(y)\,dx\,dy \notag\\
&-\sum_{\substack{n\ge2\\\log n<L}}\frac{\Lambda(n)}{\sqrt n}
\bigl(\cosh(\omega\log n)-1\bigr)
\langle f,(T_{\log n}+T_{\log n}^*)f\rangle .
\label{eq:rh-shift-form}
\end{align}
The added integral kernel is $O(|x-y|)$ at the diagonal. Together with
the finite delay sum this yields, for each fixed $L$, a constant $C_L$
such that
\begin{equation}
|Q_{\omega,L}[f]-Q_{0,L}[f]|
\leq C_L\omega^2\lVert f\rVert^2,\qquad 0\leq\omega\leq\tfrac12.
\label{eq:rh-shift-bound}
\end{equation}
The bound concerns the bounded difference of two generally unbounded
forms. Formula \eqref{eq:rh-shift-form} also specifies the endpoint
$\omega=1/2$ without dropping a distributional correction.

\paragraph{The generator, explicitly.}
Logarithmic differentiation of \eqref{eq:rh-transfer} gives the causal
generator symbol
\[
a_\omega(p)=\frac{F'(p-\omega)}{F(p-\omega)}
+\frac{F'(p+\omega)}{F(p+\omega)},\qquad
\partial_\omega K_\omega=-a_\omega K_\omega .
\]
The generator $G_{\omega,L}$ is realized from $a_\omega$ exactly as
$V_{\omega,L}$ is realized from $K_\omega$, on the same line
$\operatorname{Re}p=\eta>1$ and with the same compression:
\[
G_{\omega,L}
=M_\eta^{-1}P_L\,\mathcal M[a_\omega(\eta+i\,\cdot)]\,P_LM_\eta ,
\]
with $P_L$ multiplication by $\mathbf1_{I_L}$, $M_\eta$ multiplication by
$e^{-\eta x}$, and $\mathcal M[b]$ the whole-line Fourier multiplier by $b$.
Equivalently, $G_{\omega,L}g=(g_\omega*\widetilde g)|_{I_L}$ for the inverse
Laplace kernel $g_\omega$ of $a_\omega$, supported in $[0,\infty)$. The natural
domain is $\mathcal D_{\log,L}$ of \eqref{eq:rh-domain}, on which
$Q_{\omega,L}$ is closed and semibounded; $G_{\omega,L}$ is not symmetric, and
only its symmetric part is a form:
\[
\operatorname{Re}\langle g,G_{\omega,L}g\rangle=Q_{\omega,L}[g],
\qquad\text{equivalently}\qquad
G_{\omega,L}^*+G_{\omega,L}\ \text{represents}\ 2Q_{\omega,L}.
\]

This is the \emph{complete} shifted form, not a part of it. Writing $\xi$ as its
three factors, with $s=1/2+p$,
\[
\frac{F'}{F}(p)=\underbrace{\frac1{p+1/2}+\frac1{p-1/2}}_{\frac12s(s-1)}
\;\underbrace{-\tfrac12\log\pi+\tfrac12\psi(\tfrac14+\tfrac p2)}
 _{\pi^{-s/2}\Gamma(s/2)}
\;+\;\underbrace{\frac{\zeta'}{\zeta}(\tfrac12+p)}_{\zeta(s)},
\]
and each factor supplies one term of \eqref{eq:rh-central} and
\eqref{eq:rh-shift-form}.

\emph{The zeta factor.} Absolute convergence on $\operatorname{Re}p=\eta>1$ gives
$\frac{\zeta'}{\zeta}(\frac12+p-\omega)+\frac{\zeta'}{\zeta}(\frac12+p+\omega)
=-2\sum_{n\geq2}\Lambda(n)n^{-1/2}\cosh(\omega\log n)\,n^{-p}$, and $n^{-p}$ is
the causal delay by $\log n$. So the zeta factor contributes
$-2\sum_{\log n<L}\Lambda(n)n^{-1/2}\cosh(\omega\log n)T_{\log n}$ before
symmetrization, the cutoff being automatic because a delay of length $L$ has
zero action on $I_L$.

\emph{The gamma and $\pi$ factors.} With
$\psi(z)=-\gamma_E+\int_0^\infty(e^{-t}-e^{-zt})(1-e^{-t})^{-1}\,dt$ at
$z=\frac14+\frac{p\mp\omega}2$ and $t=2r$, these contribute a constant together
with the causal kernel $-2\cosh(\omega r)\,e^{-r/2}(1-e^{-2r})^{-1}$ at
separation $r$. At $\omega=0$ and $p=i\tau$ the multiplier is
$\operatorname{Re}\psi(\frac14+\frac{i\tau}2)-\log\pi$, the first term of
\eqref{eq:rh-central}.

\emph{The pole factors.} The four simple poles of $a_\omega$ at
$p=\pm\frac12\pm\omega$ have causal kernel
$4\cosh(\omega x)\cosh(x/2)\mathbf1_{x>0}$, which under
$\operatorname{Re}\langle g,\cdot\,g\rangle$ becomes
$2\cosh(\omega(x-y))\cosh\!\big(\frac{x-y}2\big)$; at $\omega=0$ this is
$2|C(g)|^2-2|S(g)|^2$, the contact and pole term of \eqref{eq:rh-central}.

Collecting the three, the off-diagonal kernel of $Q_{\omega,L}$ away from the
prime atoms is $\cosh(\omega(x-y))R(|x-y|)$ with $R$ as in
\eqref{eq:rh-shift-form}, and each prime atom carries
$-\Lambda(n)n^{-1/2}\cosh(\omega\log n)$. Both are the $\omega=0$ kernels
multiplied by $\cosh(\omega\cdot)$, which is \eqref{eq:rh-hyperbolic} seen on the
kernel side, and both leave the diagonal unchanged since $\cosh0=1$.

\paragraph{The pole term is invisible on the imaginary axis.}
The causal definition on $\operatorname{Re}p=\eta>1$ is not a technical
convenience here. At $p=i\tau$,
\[
\operatorname{Re}\left[\frac2{p-1/2}+\frac2{p+1/2}\right]
=\frac{2\cdot\frac12}{\frac14+\tau^2}+\frac{2\cdot(-\frac12)}{\frac14+\tau^2}=0 ,
\]
identically in $\tau$. Reading $a_\omega$ as a Fourier multiplier on the
imaginary axis therefore recovers the archimedean term and the prime delays and
\emph{loses the contact and pole term entirely}. The pole factors reach
$Q_{\omega,L}$ only through the right-half-plane Laplace realization above.

\paragraph{The gamma kinetic term as an infinitesimal shift response.}
Isolating the gamma and $\pi$ factors of $K_\omega$ gives
\[
K^\Gamma_\omega(p)
=\pi^\omega
\frac{\Gamma(1/4+p/2-\omega/2)}
     {\Gamma(1/4+p/2+\omega/2)}.
\]
Direct logarithmic differentiation yields
\[
-\left.\partial_\omega\log K^\Gamma_\omega(p)\right|_{\omega=0}
=\psi(1/4+p/2)-\log\pi.
\]
The gamma expression is regular at $p=i\tau$, so its real part is
\[
\operatorname{Re}\!\left[
-\left.\partial_\omega\log K^\Gamma_\omega(i\tau)\right|_{\omega=0}
\right]=B(\tau^2)+w_0.
\]
This is the arithmetic origin of the kinetic multiplier. The positive
tower realizes the $B$ portion of the zero-shift generator, not the
finite-shift transfer itself. Its fields therefore remain nontrivial
at $\omega=0$, although $K_0=1$ and $D_{0,L}=0$, with the contraction
defect $D_{\omega,L}$ defined below. Evaluating this gamma expression
on the imaginary axis does not replace the right-half-plane causal
definition of the full zeta transfer.

For fixed $L$, first choose a test function $f\in C_c^\infty(I_L)$:
a smooth, possibly complex-valued input compactly supported inside
$I_L$. Keep $f$ fixed as $\omega$ varies. The evolution and its
scalar energy balance are
\begin{equation}
\begin{aligned}
\frac{d}{d\omega}(V_{\omega,L}f)
&=-G_{\omega,L}(V_{\omega,L}f),\\
\frac{d}{d\omega}\lVert V_{\omega,L}f\rVert_{L^2(I_L)}^2
&=-2Q_{\omega,L}[V_{\omega,L}f].
\end{aligned}
\label{eq:rh-energy-derivative}
\end{equation}
The norm in the second line is the usual unweighted $L^2(I_L)$ norm,
so its square is $\int_{I_L}|(V_{\omega,L}f)(x)|^2\,dx$.
The brackets in $Q_{\omega,L}[V_{\omega,L}f]$ mean evaluation:
first apply the linear operator $V_{\omega,L}$ to $f$, then evaluate
the quadratic functional $Q_{\omega,L}$ on the resulting function
to obtain a scalar energy. This is composition as maps, rather than
a product of two linear operators. Products of operators below,
such as $G_{\omega,L}V_{\omega,L}$, do denote operator composition.

The first identity holds for every such $f$, so we may express it
as an \emph{operator differential equation}:
\[
\boxed{\frac{\partial V_{\omega,L}}{\partial\omega}
=-G_{\omega,L}V_{\omega,L},\qquad V_{0,L}=I.}
\]
Here $f$ has become implicit: the operator identity means that the
two sides have equal action on every admissible test function.
It describes the evolution of the entire transformation
$V_{\omega,L}$, with $V_{0,L}f=f$ as its initial condition.

The energy balance also has an operator version. The identity
\[
\lVert V_{\omega,L}f\rVert^2
=\langle f,V_{\omega,L}^*V_{\omega,L}f\rangle
\]
identifies $V_{\omega,L}^*V_{\omega,L}$ as the operator representing
the output norm squared. At positive shift, norm differentiability
as justified below permits the product rule:
\begin{align*}
\frac{\partial}{\partial\omega}(V_{\omega,L}^*V_{\omega,L})
&=\left(\frac{\partial V_{\omega,L}}{\partial\omega}\right)^*
  V_{\omega,L}
 +V_{\omega,L}^*\frac{\partial V_{\omega,L}}{\partial\omega}\\
&=-(G_{\omega,L}V_{\omega,L})^*V_{\omega,L}
  -V_{\omega,L}^*(G_{\omega,L}V_{\omega,L}).
\end{align*}
Consequently the contraction defect
$D_{\omega,L}=I-V_{\omega,L}^*V_{\omega,L}$ satisfies the operator
energy balance
\[
\boxed{\frac{\partial D_{\omega,L}}{\partial\omega}
=V_{\omega,L}^*(G_{\omega,L}^*+G_{\omega,L})V_{\omega,L},
\qquad D_{0,L}=0.}
\]
In this last formula the right-hand side is understood through
quadratic forms: $G_{\omega,L}^*+G_{\omega,L}$ represents the
symmetric generator form $2Q_{\omega,L}$, pulled back by
$V_{\omega,L}$. The preceding product-rule expression specifies the
bounded operator at positive shift without requiring the separate
application of $G_{\omega,L}^*$ to an evolved input.
Pairing with any admissible $f$ recovers the scalar energy balance
in \eqref{eq:rh-energy-derivative}; conversely, its validity for every
test function determines the symmetric form by polarization.
Thus $f$ is implicit in the operator energy balance as well.

At $\omega=0$ the evolution identity is interpreted by a right
derivative on these test functions, and $V_{0,L}=I$ is a strong
limit. No operator-norm derivative at zero is asserted.
To justify these statements, the gamma ratio
asymptotics and absolutely convergent Dirichlet series on
$p=\eta+i\tau$ give, uniformly for $0\leq\omega\leq1/2$,
\[
|\partial_\omega^jK_\omega(\eta+i\tau)|
\leq C_\eta(1+\log(2+|\tau|))^j(2+|\tau|)^{-\omega},
\qquad j=0,1,2.
\]
For positive shift these bounds give differentiability and smoothing
into $H^r$ for $0<r<\min(\omega,1/2)$, hence into the form domain.
Causality permits compression of the generator identity to $I_L$.
For smooth compactly supported initial data, rapid Fourier decay permits
differentiation at zero by dominated convergence. Integrating from a
positive shift and then taking its limit to zero gives the identities
below. No bounded inverse of $V_{\omega,L}$ is needed.

\subsection{Instantaneous energy, cumulative contraction, and the zero-shift limit}
\label{subsec:rh-contraction}

For the contraction defect just defined, integrating the energy
balance in \eqref{eq:rh-energy-derivative} on smooth compactly
supported inputs yields
\begin{align}
\langle f,D_{\omega,L}f\rangle
&=\lVert f\rVert^2-\lVert V_{\omega,L}f\rVert^2
=2\int_0^\omega Q_{s,L}[V_{s,L}f]\,ds,
\label{eq:rh-storage}\\
Q_{0,L}[f]&=
\lim_{\omega\downarrow0}\frac{\langle f,D_{\omega,L}f\rangle}{2\omega}.
\label{eq:rh-central-limit}
\end{align}
Here $s$ is a running shift. The generator controls instantaneous
change, whereas $D_{\omega,L}$ records the accumulated change along
the moving input $V_{s,L}f$.
Nonnegativity of $D_{\omega,L}$ on every input is exactly the contraction
condition $\lVert V_{\omega,L}\rVert\leq1$. Although $D_{0,L}=0$,
its initial derivative contains the nontrivial central form.

Positivity of $Q_{\omega,L}$ at all shifts is not the criterion
\eqref{eq:rh-criterion}. It is an additional requirement on the family.
Indeed, with $Xf(x)=xf(x)$, expansion of the kernels gives
\begin{equation}
Q_{\omega,L}[f]
=Q_{0,L}[\cosh(\omega X)f]-Q_{0,L}[\sinh(\omega X)f].
\label{eq:rh-hyperbolic}
\end{equation}
This is a difference of central energies; their individual
nonnegativity does not supply the comparison needed for a nonnegative
difference. Similarly, a negative instantaneous value on one input
away from zero does not determine the sign of the cumulative integral
\eqref{eq:rh-storage}.

A useful sufficient local route starts with
$Q_{0,L}\geq m_LI$, where $m_L>0$. Equations
\eqref{eq:rh-shift-bound} and \eqref{eq:rh-energy-derivative} imply
$Q_{s,L}\geq(m_L-C_Ls^2)I$, and integration gives
\begin{equation}
\lVert V_{\omega,L}\rVert
\leq\exp(-m_L\omega+C_L\omega^3/3).
\label{eq:rh-small-contraction}
\end{equation}
Thus sufficiently small positive shifts give contraction.
Their allowed range may shrink with $L$; no common positive shift
range is required for the following implication.

The central-limit identity and restriction of causal outputs to shorter
intervals give the following sufficient route directly.

\begin{proposition}[A sufficient contraction route to RH]
\label{prop:rh-diagonal}
If there exist $L_j\to\infty$ and $0<\omega_j\downarrow0$ with
$\lVert V_{\omega_j,L_j}\rVert\leq1$ for every $j$, then RH holds.
\end{proposition}
\begin{proof}
Fix $L>0$ and $f\in C_c^\infty(I_L)$. For sufficiently large $j$,
extend $f$ by zero into $I_{L_j}$ and restrict its output back to $I_L$.
The causal convolution defining the transfers gives
$\lVert V_{\omega_j,L}f\rVert\leq\lVert f\rVert$.
The quotients in \eqref{eq:rh-central-limit} are therefore nonnegative
along this sequence. Their limit gives $Q_{0,L}[f]\geq0$.
The input and length were arbitrary, so \eqref{eq:rh-criterion} applies.
\end{proof}

The shift disappears from the conclusion because it has been sent to
zero. This proposition is an auxiliary sufficient route; the central
criterion itself is the equivalence \eqref{eq:rh-criterion}.
If central positivity has already been established at arbitrarily
large lengths, the contraction argument is unnecessary.

\subsection{Positive factors as one route within the program}
\label{subsec:rh-factor}

The direct RH objective of a positive realization is an exact identity
\begin{equation}
Q_{0,L}[f]=\lVert A_{0,L}f\rVert_{\mathcal K_{0,L}}^2,
\qquad f\in C_c^\infty(I_L),
\label{eq:rh-factor-target}
\end{equation}
with $\mathcal K_{0,L}$ a positive Hilbert space and $A_{0,L}$ defined
from arithmetic data without assuming the sign of the target form.
The zero in both subscripts fixes $\omega=0$: this is the central
identity, not a claim that the construction cannot depend on a shift.
The shorter notation $A_L$ would mean $A_{0,L}$ here.
If such a construction and its matching identity hold at arbitrary
lengths, \eqref{eq:rh-criterion} proves RH directly. Separate
finite-interval factors suffice for this implication; a single
closable operator on ordinary whole-line $L^2$ is an additional
operator-realization requirement.

A shifted factorization would instead have the form
\begin{equation}
Q_{\omega,L}[f]
=\lVert A_{\omega,L}f\rVert_{\mathcal K_{\omega,L}}^2,
\qquad (\omega,L)\in\mathcal R ,
\label{eq:rh-shift-factor-target}
\end{equation}
where $\mathcal R$ is the region of parameters on which this identity
is proved, and the inputs range over the stated form domain.
Its restriction at $\omega=0$ is \eqref{eq:rh-factor-target}.
We do not require $\mathcal R$ to contain every shift at every length.
For the direct RH route, it is enough to cover $(0,L)$ at arbitrarily
large lengths. A model defined only for positive shifts would need a
proved central limit, or the contraction route in
Proposition~\ref{prop:rh-diagonal}.

Supersymmetric notation organizes a suitable closed factor into an
off-diagonal supercharge and a nonnegative partner Hamiltonian.
The arithmetic content lies in deriving \eqref{eq:rh-factor-target}.
Choosing the target's spectral square root after assuming its positivity
does not provide that derivation.

A full factor proved only on a bounded range of support lengths has
that restricted scope. A factor of $K$ alone, or an identity on a
proper subspace of inputs, also requires an additional argument to
reach the full central criterion. The gamma kinetic factor is derived
explicitly below. A full shifted factor, where available, would enter
the accumulated-energy identity at the end of the bulk discussion;
no such factor is assumed in deriving the shift-energy balance itself.

\subsection{What a bulk--boundary realization would mean}
\label{subsec:rh-bulk}

The factor $A_{\omega,L}$ need not act through a local expression in
the input coordinate alone. It could arise by solving for auxiliary
fields and evaluating their positive energy. In this interpretation,
the \emph{boundary data} are the prescribed input $f$ and the
\emph{bulk variables} are the additional degrees of freedom allowed
to adjust to that input. The construction has three steps: prescribe
the boundary input, solve for the bulk configuration of least energy,
and express that minimum using the boundary input alone. The last
step can produce a nonlocal quadratic form even when the unreduced
energy involves only local fields and derivatives.

\paragraph{What the boundary map does.}
Let $\mathcal X_{\omega,L}$ be a linear space of admissible bulk
configurations, with its specified regularity and finite-energy
conditions. A linear boundary map
\[
\operatorname{Tr}:\mathcal X_{\omega,L}\longrightarrow\mathcal B_L,
\qquad U\longmapsto\operatorname{Tr}U
\]
extracts the input visible at the boundary. Here $\mathcal B_L$ is
the space of allowed boundary functions, with a topology appropriate
to the construction; it must contain the test functions required
by the proposed matching identity. The notation denotes a
\emph{trace map}, not the matrix or operator trace.

For a geometric field, the model case is
$(\operatorname{Tr}U)(x)=U(x,0)$. The variable $x$ still runs over
the input interval: the boundary datum is the entire function
$f(x)$, not just its values at $x=\pm L/2$. If the bulk occupies
$\mathbb R\times(0,\infty)$, one can prescribe the whole lower-edge
trace to be the zero extension $\widetilde f$; we then identify
that trace with $f$ on $I_L$. For nonsmooth fields, boundary
restriction must be defined as a continuous trace in the chosen
energy-space topology, rather than by pointwise evaluation.

For an auxiliary-channel configuration
$U=(f,u_0,u_1,\ldots)$, the same notation simply means the projection
$\operatorname{Tr}U=f$. There need not be an additional continuous
coordinate or a geometric edge. In either case, the trace selects
the prescribed data; it does not determine the other components
of $U$. Their determination is the role of the energy minimization.
Conditions at other edges or at infinity belong to the bulk
domain and must also be specified.

\paragraph{Fixing the boundary and minimizing the interior.}
Suppose the system has an independently defined nonnegative
quadratic energy $\mathcal E_{\omega,L}[U]$. For each fixed $f$,
consider \emph{all} admissible configurations satisfying
$\operatorname{Tr}U=f$. Only the bulk variables vary in this
minimization; $f$, $\omega$ and $L$ remain fixed. The proposed
arithmetic matching is
\begin{equation}
Q_{\omega,L}[f]
\ \stackrel{\text{target}}{=}\
\mathcal E_{\omega,L}^{\mathrm{eff}}[f]
:=\inf_{\operatorname{Tr}U=f}\mathcal E_{\omega,L}[U].
\label{eq:rh-bulk-target}
\end{equation}
The effective energy is the least cost of sustaining that boundary
profile. A chosen extension gives an upper bound on this cost;
it gives the exact cost only when it minimizes the energy.
The admissible class must be nonempty and contain a finite-energy
configuration for every required $f$. Since every competing energy
is nonnegative, its infimum is nonnegative as well. Existence of an
actual minimizer is a further analytical question, not a consequence
of nonnegativity alone.

To see how one finds the minimizer, temporarily suppress $\omega,L$.
Write $\mathcal E[U]=e(U,U)$, where $e$ is a Hermitian
sesquilinear form. Choose any admissible lift $Rf$ with
$\operatorname{Tr}(Rf)=f$. Every competitor then has the form
\[
U=Rf+h,\qquad h\in\mathcal N:=\ker\operatorname{Tr}.
\]
Thus $\mathcal N$ consists precisely of variations that leave the
boundary fixed. If $\mathcal P f$ is a minimizer, variations by
$t h$ and $it h$, with real $t$, give the weak equilibrium equations
\[
\operatorname{Tr}(\mathcal P f)=f,\qquad
e(h,\mathcal P f)=0\quad\text{for every }h\in\mathcal N.
\]
The first equation prescribes the boundary; the second says that
the energy cannot be lowered by changing only the interior.
It is the Euler--Lagrange equation of the bulk problem. In fact,
for every $h\in\mathcal N$ it gives the exact decomposition
\[
\mathcal E[\mathcal P f+h]
=\mathcal E[\mathcal P f]+\mathcal E[h],
\]
which verifies minimality directly.

One sufficient setting for a unique linear extension
$f\mapsto\mathcal P f$ is a Hilbert bulk space with continuous trace,
a bounded linear lift $R$ from the chosen boundary space, and a
continuous form $e$ satisfying
$e(h,h)\geq c\lVert h\rVert_{\mathcal X}^2$ on $\mathcal N$ for
some $c>0$. The correction $h_f$ is then the unique solution of
$e(h,h_f)=-e(h,Rf)$ for all $h\in\mathcal N$, and
$\mathcal P f=Rf+h_f$ is linear in $f$. These are sufficient
hypotheses, not assumptions already proved for an arithmetic
completion. In degenerate or infinite-channel models, one must
instead justify the relevant completion, possible zero-energy
quotient, and attainment on the stated domain.

\paragraph{From the minimizing field to a boundary factor.}
Restore the parameters and suppose the bulk energy is given
independently as
$\mathcal E_{\omega,L}[U]
=\lVert\mathcal C_{\omega,L}U\rVert_{\mathcal K_{\omega,L}}^2$.
For a local field theory, $\mathcal C_{\omega,L}U$ can be a list
of weighted derivatives and field combinations; for a network,
it can be a list of edge differences and local amplitudes.
If a linear minimizing extension $\mathcal P_{\omega,L}$ exists,
the construction reads
\[
f\ \overset{\mathcal P_{\omega,L}}{\longmapsto}\ U_{\min}
\ \overset{\mathcal C_{\omega,L}}{\longmapsto}\ \text{energy amplitudes},
\qquad
\operatorname{Tr}\mathcal P_{\omega,L}=I.
\]
Consequently
\[
A_{\omega,L}:=\mathcal C_{\omega,L}\mathcal P_{\omega,L},
\qquad
\mathcal E_{\omega,L}^{\mathrm{eff}}[f]
=\lVert A_{\omega,L}f\rVert^2.
\]
Here the trace goes from bulk to boundary, whereas the minimizing
extension goes from boundary to bulk. Their composition above is
the identity on boundary inputs; $\mathcal P_{\omega,L}
\operatorname{Tr}$ generally replaces a bulk field by its
equilibrium extension rather than returning that field unchanged.
The factor $A_{\omega,L}$ records the energy amplitudes of this
extension. Its definition comes from solving the bulk equations,
and can therefore be nonlocal in $f$.

Polarization gives the effective sesquilinear form
$e^{\mathrm{eff}}(g,f)
=\langle A_{\omega,L}g,A_{\omega,L}f\rangle$.
If this reduced form is densely defined and closed on $L^2(I_L)$,
it determines a nonnegative self-adjoint boundary operator
$\mathsf S_{\omega,L}$, characterized by
\[
e^{\mathrm{eff}}(g,f)
=\langle g,\mathsf S_{\omega,L}f\rangle
\]
for $f$ in its operator domain and $g$ in its form domain.
When $A_{\omega,L}$ is taken as a closed factor, this is the
operator $A_{\omega,L}^*A_{\omega,L}$. The boundary map
$\operatorname{Tr}$ and the boundary operator $\mathsf S_{\omega,L}$
thus have different jobs: one selects prescribed data, while the
other describes the response after the bulk has equilibrated.
Neither the trace nor the minimizing extension is assumed bounded
in the ordinary $L^2$ norm merely because the notation resembles
that of finite matrices.

\paragraph{The same mechanism in block form.}
The algebra can be seen completely in a finite-dimensional
quadratic system. Separate the fixed boundary vector $f$ from
the free interior vector $u$, and let
\[
\mathcal E[f,u]
=
\begin{pmatrix}f\\u\end{pmatrix}^{\!*}
\begin{pmatrix}
H_{\mathrm{bb}}&H_{\mathrm{bi}}\\
H_{\mathrm{ib}}&H_{\mathrm{ii}}
\end{pmatrix}
\begin{pmatrix}f\\u\end{pmatrix},
\qquad H_{\mathrm{ib}}=H_{\mathrm{bi}}^*.
\]
Assume the full block matrix is nonnegative and
$H_{\mathrm{ii}}$ is positive definite. Here
$\operatorname{Tr}(f,u)=f$. Minimizing over $u$ gives
\[
u_{\min}=-H_{\mathrm{ii}}^{-1}H_{\mathrm{ib}}f,\qquad
\mathsf S=H_{\mathrm{bb}}
-H_{\mathrm{bi}}H_{\mathrm{ii}}^{-1}H_{\mathrm{ib}}.
\]
Indeed, completing the square yields
\[
\mathcal E[f,u]
=\langle f,\mathsf S f\rangle
+\bigl\lVert H_{\mathrm{ii}}^{1/2}
  (u+H_{\mathrm{ii}}^{-1}H_{\mathrm{ib}}f)\bigr\rVert^2.
\]
Thus $\mathcal E^{\mathrm{eff}}[f]=\langle f,\mathsf S f\rangle$;
the matrix $\mathsf S$ is the \emph{Schur complement}.
The subtracted term describes the energy reduction obtained by
allowing the interior to relax. Positivity of $\mathsf S$ follows
from positivity of the \emph{full} energy matrix. Positivity of the
interior block alone would not suffice. Even if the original
couplings are local or sparse, the inverse interior operator can
couple distant boundary inputs. This is the algebraic reason that
eliminating local bulk fields can produce a nonlocal boundary form.
For differential operators or infinitely many channels, the same
calculation requires a form-domain and convergence argument.

In a geometric extension this reduced response is often called a
\emph{Dirichlet-to-Neumann map}: prescribe the boundary value
(Dirichlet data), solve the equilibrium equation, and read the
outward normal flux (Neumann data). For a regular weighted energy
$\int a(y)|\nabla U|^2\,dx\,dy$, integration by parts in an
equilibrium field identifies the effective energy with the
boundary pairing against that flux, with the remaining boundary
terms fixed by the domain. The trace itself only reads $U(x,0)$;
the flux response also depends on the solved field.
Positive extension realizations of functions of the Laplace
operator provide a rigorous setting for this mechanism
\cite{KwasnickiMuchaBackground}.

\paragraph{What must match the arithmetic form.}
The qualifier \emph{target} in \eqref{eq:rh-bulk-target} separates
the variational construction from the arithmetic theorem still
required. One must specify a bulk energy and its domains from
independent data, solve or characterize the constrained minimum,
and compute the resulting boundary form. That computation must
reproduce $Q_{\omega,L}$ with its exact normalization and signs.
For RH, the required central identity includes the gamma term,
the pole terms and every active prime delay at arbitrarily large
lengths. A successful match then supplies
$Q_{0,L}[f]=\lVert A_{0,L}f\rVert^2$ and invokes
\eqref{eq:rh-criterion}. A positive bulk system by itself proves
positivity of its own reduced form; identifying that form with the
Weil form is the substantive research problem.

\paragraph{An explicit positive realization of the gamma kinetic term.}
We now derive the bulk mechanism for the kinetic component of the
opening decomposition. The construction uses the positive partial
fractions of the digamma function and does not assume the sign of
the full Weil form. Recall $a_k=2k+1/2$ and
\begin{align}
B(s)&=\sum_{k=0}^{\infty}\frac2{a_k}\frac{s}{a_k^2+s},
\qquad s\geq0,\notag\\
K[\widetilde f]&=
\frac1{2\pi}\int_{\mathbb R}B(\tau^2)|\widehat f(\tau)|^2\,d\tau .
\label{eq:rh-kinetic}
\end{align}
The partial-fraction calculation in the opening subsection identifies
this series with the gamma multiplier. Equivalently, set
\[
j(y)=\sum_{k=0}^\infty e^{-a_ky}
=\frac{e^{-y/2}}{1-e^{-2y}},\qquad y>0.
\]
Then
\[
\begin{aligned}
2\int_0^\infty e^{-ay}(1-\cos(\tau y))\,dy
&=\frac2a\frac{\tau^2}{a^2+\tau^2},\\
B(\tau^2)&=2\int_0^\infty j(y)(1-\cos(\tau y))\,dy.
\end{aligned}
\]
Tonelli's theorem applies to these nonnegative integrands. Near zero,
$j(y)\sim1/(2y)$ and $1-\cos(\tau y)=O(y^2)$; at infinity $j$
decays exponentially. This verifies convergence of the integral.

The term \emph{kinetic} describes the dependence on spatial variation.
Here the response is nonlocal, with logarithmic growth at high
frequency. It is not the ordinary local $\lVert f'\rVert^2$ energy.
For example, writing $F=\widetilde f$ and applying Plancherel,
\[
K[F]=\int_0^\infty j(y)
\lVert F-F(\,\cdot-y)\rVert_{L^2(\mathbb R)}^2\,dy.
\]
The fields below realize this response by energies local in $x$
before minimization.

For auxiliary fields $u_k\in H^1(\mathbb R)$, define
\begin{equation}
\mathcal E^{\mathrm{kin}}[\widetilde f,(u_k)]
=\sum_{k=0}^{\infty}\frac2{a_k}
\int_{\mathbb R}
\left(|\widetilde f-u_k|^2+a_k^{-2}|u_k'|^2\right)\,dx .
\label{eq:rh-kinetic-bulk}
\end{equation}
Every term is nonnegative and local in $x$ before elimination.
In the general notation above,
$\operatorname{Tr}(\widetilde f,(u_k))=f$: the first component is
fixed and the fields $u_k$ are varied independently.
Their weak equilibrium equations are
\[
(1-a_k^{-2}\partial_x^2)u_k=\widetilde f
\quad\text{on }\mathbb R.
\]
Solving these equations gives the minimizing extension
$\mathcal P f=(\widetilde f,(u_k))$.
Completion of the square at each Fourier frequency gives
\begin{equation}
\inf_{(u_k)}\mathcal E^{\mathrm{kin}}[\widetilde f,(u_k)]
=K[\widetilde f],\qquad
\widehat u_k(\tau)=\frac{a_k^2}{a_k^2+\tau^2}\widehat f(\tau).
\label{eq:rh-kinetic-reduction}
\end{equation}
Here is the full one-channel calculation. At frequency $\tau$, put
$\lambda=\tau^2/a^2$, $h=\widehat f(\tau)$ and $v=\widehat u(\tau)$.
Completing the square gives
\[
|h-v|^2+\lambda|v|^2
=(1+\lambda)\left|v-\frac{h}{1+\lambda}\right|^2
+\frac{\lambda}{1+\lambda}|h|^2.
\]
The unique minimizer and its real-space expression are
\[
\widehat u(\tau)=\frac{a^2}{a^2+\tau^2}\widehat f(\tau),
\qquad
u(x)=\frac a2\int_{\mathbb R}e^{-a|x-y|}\widetilde f(y)\,dy.
\]
For every $f\in L^2(I_L)$ this solution belongs to $H^2(\mathbb R)$,
so it is an admissible $H^1$ field and solves the equilibrium equation.
Its smoothing length is $a^{-1}$. Multiplying the minimum by $2/a$
and applying Plancherel gives
\[
\inf_{u\in H^1(\mathbb R)}
\frac2a\int_{\mathbb R}
\left(|\widetilde f-u|^2+a^{-2}|u'|^2\right)\,dx
=\frac1{2\pi}\int_{\mathbb R}
\frac2a\frac{\tau^2}{a^2+\tau^2}|\widehat f(\tau)|^2\,d\tau.
\]
Every tower configuration has energy at least the sum of these channel
minima. The family of channel minimizers attains that sum whenever it
is finite. Monotone convergence of the nonnegative integrands gives
the kinetic reduction above; the infimum and sum have not been
interchanged without justification. Since
$B(\tau^2)=\log(2+|\tau|)+O(1)$, finite minimized energy is equivalent
to $f\in\mathcal D_{\log,L}$. The fields extend over the whole line
even when the input is supported in $I_L$, and their exterior tails
are included. Truncating the fields at the input endpoints would
change the reduced energy.

The corresponding bulk amplitude map is explicitly
\[
\mathcal C(\widetilde f,(u_k))
=\left(
\sqrt{\frac2{a_k}}(\widetilde f-u_k),
\sqrt{\frac2{a_k^3}}\,u_k'
\right)_{k\geq0}.
\]
Its target is the Hilbert direct sum of the displayed $L^2(\mathbb R)$
components. Thus the factor produced by elimination is
$A^{\mathrm{kin}}f=\mathcal C\mathcal P f$, and
$\lVert A^{\mathrm{kin}}f\rVert^2=K[\widetilde f]$.
The minimizing channel configuration has finite total energy precisely
when the displayed kinetic integral is finite. This identifies the
trace, equilibrium extension, amplitude map and reduced energy in
one concrete instance of the general construction.

\paragraph{What this construction establishes for the program.}
The calculation realizes the arithmetic kinetic component as the
minimum of an independently nonnegative local field energy:
\[
\boxed{
K[\widetilde f]
=\inf_{(u_k)}\mathcal E^{\mathrm{kin}}[\widetilde f,(u_k)]
=\lVert A^{\mathrm{kin}}f\rVert^2\geq0.}
\]
This is the first constructive success of the proposed mechanism:
one prescribed component is obtained exactly, with explicit field
equations, an attained minimum and the required input domain. It
also supplies a positive system whose couplings can be modified to
seek the remaining arithmetic contributions. The masses and weights
are chosen from the gamma identity; their realization does not show
that supersymmetry selects them.

The full target is the boxed decomposition in the opening subsection.
The full central Weil form additionally contains
$w_0\lVert f\rVert^2$, where $w_0=\psi(1/4)-\log\pi<0$,
the pole terms, and the prime sum in \eqref{eq:rh-central}.
Explaining these contributions within a common positive bulk energy
is the remaining completion problem. The kinetic construction does
not by itself solve it.

\paragraph{Where the shift can enter a bulk description.}
There are two distinct uses of $\omega$. In
\eqref{eq:rh-bulk-target}, it is a parameter of the energy and boundary
problem; the direct RH target evaluates that family at $\omega=0$.
Alternatively, if the full generator factors are available for
$0\leq s\leq\omega$, the shift-energy identity becomes
\begin{equation}
\langle f,D_{\omega,L}f\rangle
=2\int_0^\omega
\lVert A_{s,L}V_{s,L}f\rVert_{\mathcal K_{s,L}}^2\,ds .
\label{eq:rh-shift-bulk-energy}
\end{equation}
This describes an accumulated positive energy on the strip
$0<s<\omega$, with initial data $f$ and trajectory
$g(s)=V_{s,L}f$. Its boundary observable is the contraction defect
$D_{\omega,L}$; the central form is recovered from its initial rate
by \eqref{eq:rh-central-limit}. This is a conditional identity: the full factors must first exist
throughout the integration interval on the evolved inputs, with
measurable and integrable squared amplitudes. The kinetic factor
$A^{\mathrm{kin}}$ alone cannot be substituted for a full generator
factor. An independent bulk explanation would also have to derive
the factor and evolution
from the specified system; writing the integral after those objects
are known does not supply the missing arithmetic mechanism.

The coordinate $x$ labels the original input, $L$ limits its support,
and a geometric coordinate $y$ or channel index $k$ labels additional
degrees of freedom. The zeta shift $\omega$ is not automatically that
geometric coordinate. Using the running shift $s$ as an evolution
coordinate in \eqref{eq:rh-shift-bulk-energy} is a specific
interpretation with the displayed energy identity as its content.
No AdS geometry, conformal field theory, or interacting supersymmetric
field theory is assumed or established by these constructions.

\begin{thebibliography}{9}
\bibitem{SuzukiScrewBackground}
M. Suzuki, \emph{Weil's quadratic form via the screw function}, 2026,
version 2, especially \S1.1.
\href{https://arxiv.org/abs/2606.09096v2}{arXiv:2606.09096v2}.
\bibitem{DLMF}
NIST Digital Library of Mathematical Functions, Chapter 5,
especially \S\S5.4, 5.7, 5.9 and 5.11.
\href{https://dlmf.nist.gov/5}{DLMF, Chapter 5}.
\bibitem{KwasnickiMuchaBackground}
M. Kwa\'snicki and J. Mucha,
\emph{Extension technique for complete Bernstein functions of the Laplace operator},
Journal of Evolution Equations \textbf{18} (2018), 1341--1379.
\href{https://doi.org/10.1007/s00028-018-0444-4}{doi:10.1007/s00028-018-0444-4};
\href{https://arxiv.org/abs/1707.02475}{arXiv:1707.02475},
especially the variational formulation and its measure-coefficient extension.
\end{thebibliography}
\end{document}
